\documentclass[11pt]{article}
\usepackage{amsmath,amssymb,amsthm}
\usepackage[margin=1in]{geometry}
\newtheorem{theorem}{Theorem}
\newtheorem{lemma}{Lemma}
\begin{document}
\title{V107: Exact-Stretch Avoidance for Essential Ternary Signed Majority}
\author{NC0\_k-Avoid Laboratory}
\date{}
\maketitle

\paragraph{Setting.}
Every output gate is
\[
 \operatorname{MAJ}(x_a\oplus p_a,x_b\oplus p_b,x_c\oplus p_c)
\]
on three distinct inputs.  Its three pair transports are
$d_{ij}=1\oplus p_i\oplus p_j$ and satisfy
$d_{01}\oplus d_{02}\oplus d_{12}=1$.

\begin{lemma}[Minimal surplus]
Every output family of positive support surplus contains an inclusion-minimal
subfamily $F$ with
\[
 |F|=|N(F)|+1,
\]
and every proper $J\subset F$ satisfies $|N(J)|\ge |J|$.
\end{lemma}

\begin{lemma}[Frame transversal after one deletion]
Delete any gate $g\in F$.  There is a polynomial-time constructible choice of
one pair candidate from every gate of $F\setminus\{g\}$ that is independent in
the signed-frame matroid.
\end{lemma}

\begin{proof}
For every gate subfamily $J$, the union of its three pair candidates has frame
rank $|N(J)|$, because every connected component contains an unbalanced majority
triangle.  The minimal-surplus Hall inequalities therefore satisfy Rado's
matroid-marriage condition.  Use ordinary polynomial-time matroid intersection
between the partition matroid and the signed-frame matroid.
\end{proof}

\begin{lemma}[Rank-tight shape]
The selected transversal has $r=|N(F)|$ edges on $r$ vertices.  Every connected
component is an unbalanced unicyclic graph.
\end{lemma}

\begin{lemma}[Three-terminal compression]
If the three variables of the omitted gate lie in one selected component, retain
the unique unbalanced cycle and the minimal attachment forest for the three
terminals, then suppress every nonterminal degree-two vertex.  The resulting
signed virtual kernel has at most six vertices and at most six virtual paths.
Each virtual path implication lifts to an implication path in the original
selected gates by XOR-composing its transport labels.
\end{lemma}

\begin{proof}
If $k$ terminals lie off the cycle in $c$ nonempty attachment trees, reduced
attachment trees contribute at most $2k-c$ paths.  The suppressed cycle has at
most $c+(3-k)$ marked arcs.  Thus the total is at most
\[
 (2k-c)+(c+3-k)=k+3\le6.
\]
Sequentially targeting the actual gates of a suppressed path realizes either
chosen source phase and transports it by the XOR of the path labels, together
with the contrapositive implication.
\end{proof}

\begin{lemma}[Computer-assisted kernel lemma]
Every reduced three-terminal unbalanced-unicyclic kernel allowed by the previous
lemma, for every signed edge-parity assignment and every one of the eight
signed-majority polarities on the terminals, admits path phases and one omitted
gate target whose virtual 2-CNF formula is unsatisfiable.
\end{lemma}

\begin{proof}[Finite exhaustive verification]
A generator enumerates the one-, two-, and three-root cycle cases, all
surjective assignments of the three labeled terminals to roots, all choices of
terminals lying directly on roots, and all reduced rooted attachment trees via
Pruefer sequences with the permitted degree-at-least-three Steiner vertices.
It produces 164 kernel descriptions.  Enumerating every virtual-edge parity
assignment with odd unique-cycle parity and all eight missing-gate polarities
gives 16,032 exact cases.  Each case has at most six virtual paths, so at most
$2^6\cdot2=128$ phase/target choices are exhausted.  Every case has an
unsatisfiable choice.  A separate direct census on unreduced simple unicyclic
graphs provides additional coverage controls.
\end{proof}

\begin{theorem}[Essential signed-majority exact-stretch avoidance]
Let
\[
 C:\{0,1\}^n\to\{0,1\}^m,\qquad m>n,
\]
be an explicit circuit whose outputs are essential ternary signed-majority
gates.  A word outside $\operatorname{Range}(C)$ can be constructed
deterministically in polynomial time.
\end{theorem}

\begin{proof}
Extract a minimal-surplus block and omit one gate.  Construct the frame-independent
pair transversal.  If the omitted support meets multiple components, one of its
pair clauses joins two unbalanced cycles and yields the V105 loose-handcuff
contradiction.  Otherwise compress the common component to the constant kernel,
find an unsatisfiable phase/target assignment by the finite lemma, and lift the
virtual clauses to exact original majority clauses.  Extra clauses supplied by
each fixed majority output only strengthen the resulting 2-CNF formula.  Thus
the chosen block output has no preimage, and arbitrary values outside the block
extend it to a missing full output.  All nonconstant steps are polynomial and
the kernel search has constant size.
\end{proof}

\paragraph{Boundary.}
This theorem concerns only the essential ternary signed-majority predicate
orbit.  It does not solve the unrestricted MUX/bijunctive orbit, general
$NC^0_3$-Avoid, or P versus NP.  Novelty and priority are not established.

\end{document}
